\documentclass[12pt]{article}
\usepackage[utf8]{inputenc}
\usepackage[T1]{fontenc}
\usepackage{hyperref}
\title{Wirangu Language Resources for K-12 Curriculum Development}
\author{Compiled by Claude Sonnet 5.5 (claude-sonnet-5-5) for Andr\'as Kornai}
\date{}
\begin{document}
\maketitle

\section{Introduction}
Wirangu (ISO 639-3 wgu; Glottocode wira1265) is a Pama-Nyungan language of the Ceduna region on the west coast of South Australia. Wikipedia describes it as moribund, with two semi-speakers in 2007. The sources found are mostly older, so current status should be checked with the community.

\section{Primary Educational Resources}
Linguist Paul Monaghan worked with elder Gladys Miller on a talking picture dictionary of about 200 common Wirangu words, on compact disc with pronunciation and sentences. An illustrated children's book in Wirangu about a wombat hunt was also launched. Both are described in a 2006 University of Adelaide news item: \url{https://www.adelaide.edu.au/adelaidean/issues/12741/news12757.html}. Whether they are still available needs to be confirmed.

The Wirangu Language Project archive at the Aboriginal \&amp; Torres Strait Islander Data Archive holds a major dictionary and materials produced for language revival in schools and the community: \url{https://www.atsida.edu.au/archive/datasets/wlp}. Access is mixed (some open, some conditional, some closed), and some material is restricted to men. See also \url{https://www.atsida.edu.au/node/645}.

\section{Linguistic and Reference Materials}
The ATSIDA archive above holds the dictionary. A general overview is at \url{https://en.wikipedia.org/wiki/Wirangu_language}.

\section{Religious and Community Resources}
The Far West Languages Centre, a program of the Ceduna Aboriginal Corporation, supports Wirangu, Gugada and Mirning: \url{https://fwlc.org.au/about-1}. It is likely the best contact point for current materials. The Mobile Language Team has a trainee-led Wirangu place names project for the Ceduna area: \url{https://mobilelanguageteam.com.au/new-aboriginal-language-trainee-at-the-mlt/}. No scripture materials were found.

\section{Digital and Audio Resources}
The National Film and Sound Archive curates a film about elders Doreen and Gladys Miller and the community effort to write books to teach Wirangu to younger people: \url{https://www.nfsa.gov.au/collection/curated/wirrangul-women-always-have-always-will-wirangu-language}.

\section{Curriculum Development Context}
Wirangu is not known to be taught in schools. The ATSIDA dictionary and project materials, the talking dictionary, the children's book, and the Far West Languages Centre make up a small base for a heritage and revitalization curriculum. Materials should be used only with the permission of the community. Restrictions on some archive items must be respected.
\end{document}